%% kapa: exam=1 round=28 year=2017 period=1 form=A subject=civil_law subject_ko=민법 seq=39 number=39
%% source: Q-Net 감정평가사 기출문제 (cst00309) artlSeq=5199603 file=extracted/1008/5199603/2017년 제28회 감정평가사 1차시험 문제지/1교시 A형.pdf
%% conversion: kordoc 4.24.51 -> tools/gen1.py (auto); review=proofread-sonnet; flags=none
\begin{question}{39}
가등기담보 등에 관한 법률상 가등기담보에 대한 설명으로 옳은 것은? (다툼이 있으면 판례에 따름)
\begin{choices}
  \item 후순위권리자는 청산기간 동안에는 담보목적부동산의 경매를 청구할 수 없다.
  \item 채무자는 청산기간이 지나기 전이라도 후순위권리자에 대한 통지 후 청산금에 관한 권리를 제3자에게 양도하면 이로써 후순위권리자에게 대항할 수 있다.
  \item 담보목적물에 대한 사용ㆍ수익권은 채무자에게 지급되어야 할 청산금이 있더라도 그 지급없이 청산기간이 지나면 채권자에게 귀속된다.
  \item 담보가등기를 마친 부동산이 강제경매를 통해 매각되어도, 담보가등기권리는 피담보채권액 전부를 변제받지 않으면 소멸하지 않는다.
  \item 담보가등기를 마친 부동산에 대하여 강제경매가 개시된 경우, 담보가등기를 마친 때를 기준으로 담보가등기권리자의 순위가 결정된다.
\end{choices}
\end{question}
\answer{5}
